\subsection{THE CARDINAL OF A SET}

DEFINITION 1. A set X is said to be equipotent to a set Y if there exists a bijection of X onto Y. The relation ``X is equipotent to Y'' is denoted by Eq(X, Y).

The relations Eq(X, Y) and Eq(Y, X) are clearly equivalent, so that the relation Eq(X, Y) is symmetric; when it is true, we say that X and Y are equipotent. Next, Eq(X, X) is true. Finally, the relation Eq(X, Y) is transitive since the composition of two bijections is a bijection (Chapter II, §3, no. 8, Theorem 1); it is therefore an equivalence relation, reflexive on every set.

From what has been said it follows that if X and Y are equipotent, the relation

\[
(\forall \mathbf {Z}) (\operatorname{Eq} (\mathbf {X}, \mathbf {Z}) \Longleftrightarrow \operatorname{Eq} (\mathbf {Y}, \mathbf {Z}))
\]

is true. Now the axiom scheme S7 (Chapter I, §5, no. 1) gives us the following axiom :

\[
((\forall Z) (\operatorname{Eq} (X, Z) \Longleftrightarrow \operatorname{Eq} (Y, Z)) \Longrightarrow (\tau_ {Z} (\operatorname{Eq} (X, Z)) = \tau_ {Z} (\operatorname{Eq} (Y, Z))).
\]

Hence, if X and Y are equipotent, we have

\[
\tau_ {z} (\operatorname{Eq} (X, Z)) = \tau_ {z} (\operatorname{Eq} (Y, Z)),
\]

which justifies the following definition :

DEFINITION 2. The set \(\tau_{\mathbf{Z}}(\operatorname{Eq}(\mathbf{X},\mathbf{Z}))\) is called the cardinal of \(\mathbf{X}\) (or the power of \(\mathbf{X}\) ) and is written \(\operatorname{Card}(\mathbf{X})\) .

Since Eq(X, X) is true, Card(X) is equipotent to X (Chapter I, §4, Scheme S5). We have therefore proved the following result:

PROPOSITION 1. Two sets X and Y are equipotent if and only if their cardinals are equal.

\minorheading{Examples}

\begin{enumerate}
\def\labelenumi{(\arabic{enumi})}
\item
  \(\operatorname{Card}(\phi)\) is denoted by 0. Since the only set equipotent to \(\phi\) is \(\phi\) (Chapter II, §3, nos. 1 and 4), we have \(0 = \operatorname{Card}(\phi) = \phi\) .
\item
  All sets consisting of a single element are equipotent, since \(\{(a, b)\}\) is the graph of a bijection of \(\{a\}\) onto \(\{b\}\) ; in particular, they are all equipotent to \(\{\varnothing\}\) . The cardinal \(\operatorname{Card}(\{\varnothing\}) = \tau_{\mathbf{Z}}(\operatorname{Eq}(\{\varnothing\}, \mathbf{Z}))\) is denoted by 1. (*)
\item
  \(\operatorname{Card}(\{\varnothing, \{\varnothing\}\})\) is denoted by 2; this is the cardinal of every set consisting of two distinct elements.
\end{enumerate}

*(4) A Hilbert space of countable type is equipotent to the set of real numbers. *

